% GENERATED FILE. Edit canonical lessons, then run npm run generate:books.
% canonical-source-hash: fnv1a64:f7bb979901303c75
% canonical-lessons: TA-C232-palvali, TA-C232-kayam, TA-C232-mattirai, TA-C232-noyali, TA-C232-ceviliyar

\chapter{Toothache, Wound, Tablet, Patient, Nurse}
\label{ch:ta-toothache-wound-tablet-patient-nurse}
\hlchaptermodality{232}

\begin{chapteropening}
\textbf{By the end of this chapter:} \emph{I can say a toothache, a wound, a tablet, a patient and a nurse.}

Complete the last lesson of chapter 232: I can say a toothache, a wound, a tablet, a patient and a nurse.
\end{chapteropening}

\section[\texorpdfstring{\ta{பல்வலி} (palvali)}{பல்வலி (palvali)}]{\ta{பல்வலி} (palvali) — a toothache}

\label{lesson:TA-C232-palvali}



\begin{quote}
Before the new one: say the Tamil for to beg, then the Tamil for to manage.
\end{quote}

\subsection*{You'll want to know: \ta{பல்வலி}}
\textbf{\ta{பல்வலி}} — \emph{palvali} — \textquotedblleft{}a toothache\textquotedblright{}.

\begin{grammarlens}[title={What you've built}]
One more word for this chapter.
\end{grammarlens}

\subsection*{Your turn}
\begin{itemize}
\raggedright
  \item \emph{Say it:} \emph{palvali}
  \item \emph{Say it:} \emph{palvali}, once more
  \item \emph{Say it:} say \emph{palvali}
  \item \emph{Recall:} say the Tamil for to beg, then the Tamil for to manage, then say \emph{palvali} again
\end{itemize}

\begin{tcolorbox}[breakable,colback=teal!4,colframe=teal!35!black,title={Before you move on}]
What does \ta{பல்வலி} mean? (\textquotedblleft{}A toothache\textquotedblright{}.) Say it once more.
\end{tcolorbox}
\section[\texorpdfstring{\ta{காயம்} (kāyam)}{காயம் (kāyam)}]{\ta{காயம்} (kāyam) — a wound, an injury}

\label{lesson:TA-C232-kayam}



\begin{quote}
Before the new one: say the Tamil for to manage, then the Tamil for a toothache.
\end{quote}

\subsection*{You'll want to know: \ta{காயம்}}
\textbf{\ta{காயம்}} — \emph{kāyam} — \textquotedblleft{}a wound, an injury\textquotedblright{}.

\begin{grammarlens}[title={What you've built}]
One more word for this chapter.
\end{grammarlens}

\subsection*{Your turn}
\begin{itemize}
\raggedright
  \item \emph{Say it:} \emph{kāyam}
  \item \emph{Say it:} \emph{kāyam}, once more
  \item \emph{Say it:} say \emph{kāyam}
  \item \emph{Recall:} say the Tamil for to manage, then the Tamil for a toothache, then say \emph{kāyam} again
\end{itemize}

\begin{tcolorbox}[breakable,colback=teal!4,colframe=teal!35!black,title={Before you move on}]
What does \ta{காயம்} mean? (\textquotedblleft{}A wound, an injury\textquotedblright{}.) Say it once more.
\end{tcolorbox}
\section[\texorpdfstring{\ta{மாத்திரை} (māttirai)}{மாத்திரை (māttirai)}]{\ta{மாத்திரை} (māttirai) — a tablet, a pill}

\label{lesson:TA-C232-mattirai}



\begin{quote}
Before the new one: say the Tamil for a toothache, then the Tamil for a wound.
\end{quote}

\subsection*{You'll want to know: \ta{மாத்திரை}}
\textbf{\ta{மாத்திரை}} — \emph{māttirai} — \textquotedblleft{}a tablet, a pill\textquotedblright{}.

\begin{grammarlens}[title={What you've built}]
One more word for this chapter.
\end{grammarlens}

\subsection*{Your turn}
\begin{itemize}
\raggedright
  \item \emph{Say it:} \emph{māttirai}
  \item \emph{Say it:} \emph{māttirai}, once more
  \item \emph{Say it:} say \emph{māttirai}
  \item \emph{Recall:} say the Tamil for a toothache, then the Tamil for a wound, then say \emph{māttirai} again
\end{itemize}

\begin{tcolorbox}[breakable,colback=teal!4,colframe=teal!35!black,title={Before you move on}]
What does \ta{மாத்திரை} mean? (\textquotedblleft{}A tablet, a pill\textquotedblright{}.) Say it once more.
\end{tcolorbox}
\section[\texorpdfstring{\ta{நோயாளி} (nōyāḷi)}{நோயாளி (nōyāḷi)}]{\ta{நோயாளி} (nōyāḷi) — a patient}

\label{lesson:TA-C232-noyali}



\begin{quote}
Before the new one: say the Tamil for a wound, then the Tamil for a tablet.
\end{quote}

\subsection*{You'll want to know: \ta{நோயாளி}}
\textbf{\ta{நோயாளி}} — \emph{nōyāḷi} — \textquotedblleft{}a patient\textquotedblright{}.

\begin{grammarlens}[title={What you've built}]
One more word for this chapter.
\end{grammarlens}

\subsection*{Your turn}
\begin{itemize}
\raggedright
  \item \emph{Say it:} \emph{nōyāḷi}
  \item \emph{Say it:} \emph{nōyāḷi}, once more
  \item \emph{Say it:} say \emph{nōyāḷi}
  \item \emph{Recall:} say the Tamil for a wound, then the Tamil for a tablet, then say \emph{nōyāḷi} again
\end{itemize}

\begin{tcolorbox}[breakable,colback=teal!4,colframe=teal!35!black,title={Before you move on}]
What does \ta{நோயாளி} mean? (\textquotedblleft{}A patient\textquotedblright{}.) Say it once more.
\end{tcolorbox}
\section[\texorpdfstring{\ta{செவிலியர்} (ceviliyar)}{செவிலியர் (ceviliyar)}]{\ta{செவிலியர்} (ceviliyar) — a nurse}

\label{lesson:TA-C232-ceviliyar}



\begin{quote}
Before the new one: say the Tamil for a tablet, then the Tamil for a patient.
\end{quote}

\subsection*{You'll want to know: \ta{செவிலியர்}}
\textbf{\ta{செவிலியர்}} — \emph{ceviliyar} — \textquotedblleft{}a nurse\textquotedblright{}.

\begin{grammarlens}[title={What you've built}]
One more word for this chapter.
\end{grammarlens}

\subsection*{Your turn}
\begin{itemize}
\raggedright
  \item \emph{Say it:} \emph{ceviliyar}
  \item \emph{Say it:} \emph{ceviliyar}, once more
  \item \emph{Say it:} say \emph{ceviliyar}
  \item \emph{Recall:} say the Tamil for a tablet, then the Tamil for a patient, then say \emph{ceviliyar} again
\end{itemize}

\begin{tcolorbox}[breakable,colback=teal!4,colframe=teal!35!black,title={Before you move on}]
What does \ta{செவிலியர்} mean? (\textquotedblleft{}A nurse\textquotedblright{}.) Say it once more.
\end{tcolorbox}
